\begin{apartats}

\apartat{0,75}
Determina el domini i estudia la continuïtat de les funcions següents:
\begin{graella}{2}
  \sa y=\sqrt{x^2-9} & \sa y=\ln(4-x)
\end{graella}

\begin{solucio}
i) Cal $x^2-9\ge0$, és a dir $|x|\ge3$: $\mathrm{Dom}=(-\infty,-3]\cup[3,+\infty)$.
És contínua a tot el domini (composició de contínues).\\
ii) Cal $4-x>0$: $\mathrm{Dom}=(-\infty,4)$. És contínua a tot el domini.
\end{solucio}

\apartat{1}
Troba els punts en què la funció
\[
f(x)=\frac{x^2-x-6}{x^2-2x-3}
\]
és discontínua i classifica'n la discontinuïtat.

\begin{solucio}
$f(x)=\dfrac{(x-3)(x+2)}{(x-3)(x+1)}=\dfrac{x+2}{x+1}$ per a $x\ne3$.
El domini és $\mathbb{R}\setminus\{-1,3\}$.\\
$x=3$: $\lim_{x\to3}f(x)=\dfrac54$ existeix però $f(3)$ no. \textbf{Discontinuïtat evitable.}\\
$x=-1$: els laterals valen $-\infty$ i $+\infty$. \textbf{Salt infinit} (asímptota vertical $x=-1$).
\end{solucio}

\apartat{0,75}
\textbf{Inventa.} Escriu una funció racional $g$ que compleixi simultàniament les
condicions següents: $\lim_{x\to+\infty}g(x)=3$, presenta una discontinuïtat de salt
infinit en $x=2$ i una discontinuïtat evitable en $x=-1$.

\begin{solucio}
Per exemple $g(x)=\dfrac{3(x+1)(x-5)}{(x+1)(x-2)}$.\\
El factor $(x+1)$ es cancel·la: discontinuïtat evitable en $x=-1$, amb límit
$\frac{3(-6)}{-3}=6$. En $x=2$ el denominador s'anul·la i el numerador no: salt infinit.
Numerador i denominador tenen el mateix grau i el quocient dels coeficients principals
és $3$, així que $\lim_{x\to+\infty}g(x)=3$. (La resposta no és única.)
\end{solucio}

\end{apartats}
